% FY27 ONR YIP informal brief topic description. Compile using pdflatex.
% Scientific framing tailored for morrison. Morrison PO line deliberately configurable.
\documentclass[11pt,letterpaper]{article}
\usepackage[letterpaper,top=0.60in,bottom=0.58in,left=0.72in,right=0.72in]{geometry}
\usepackage{mathptmx}
\usepackage[T1]{fontenc}
\usepackage{microtype}
\usepackage{enumitem}
\usepackage{xcolor}
\usepackage[colorlinks=true,citecolor=blue,linkcolor=black,urlcolor=blue]{hyperref}
\definecolor{accent}{HTML}{233D57}
\pagestyle{empty}
\setlength{\parindent}{0pt}
\setlength{\parskip}{3pt}
\setlist[itemize]{leftmargin=1.13em,topsep=2pt,itemsep=2pt,parsep=0pt}
\emergencystretch=1.5em
\newcommand{\head}[1]{\textbf{#1}\enspace}
\newcommand{\POline}{Dr. Jeffrey Morrison, Code 341 (assignment unconfirmed) --- Command Decision Making}
\newcommand{\FocusArea}{Artificial Intelligence and Autonomy --- Intelligent and Collaborative Agents}
\begin{document}
\begin{center}
{\fontsize{14}{16}\selectfont\bfseries Knowledge-Grounded AI Reasoning for Dynamic Decision Support\par}
\vspace{2pt}
{\small Brief topic description for Program Officer discussion --- ONR FY27 YIP (N0001426SF004)}
\end{center}
\vspace{-2pt}
{\small \textbf{Xinya Du}, Assistant Professor, Computer Science, The University of Texas at Dallas \quad \href{mailto:xinya.du@utdallas.edu}{xinya.du@utdallas.edu} \quad \href{https://xinyadu.github.io}{https://xinyadu.github.io}\par
\textbf{Program Officer:} \POline\par
\textbf{Proposed ONR S\&T Focus Area:} \FocusArea\par}
\vspace{2pt}\hrule\vspace{5pt}
\head{Problem.} Decision makers must combine reports describing evolving events, often from incomplete, contradictory, or out-of-sequence sources. One report may say that a port has closed, while a later message corrects its reopening time; a related message may describe why traffic is still delayed. An AI assistant that simply summarizes the latest retrieved passages can miss event relationships or present a disputed conclusion as certain. The important question is not only how to represent events, but how to use updated knowledge to answer decision-relevant questions.
\par\head{Central hypothesis.} Agents that maintain event-centered knowledge with uncertainty and provenance, and actively identify which new information changes a decision, can provide more reliable answers than retrieval-only summaries. We will study evidence-grounded reasoning and proactive information synthesis, including distributed reports when messages arrive in different orders.
\par\head{Research plan.}
\begin{itemize}
\item \textbf{Maintain evolving event knowledge.} Extract actors, events, times, causes, and sources from long documents and incoming updates. Preserve alternative accounts and temporal corrections instead of collapsing disagreements into a single narrative. Event-graph extraction supplies the foundation, while deciding how an update changes downstream knowledge is a new target.
\item \textbf{Reason over changes and competing explanations.} Develop multi-hop reasoning that traces a question to the relevant events and evidence, checks temporal and causal consistency, and identifies which assumptions remain unresolved. We will test whether event-centered reasoning yields more accurate answers and clearer uncertainty than flat document retrieval.
\item \textbf{Proactively select and share decision-relevant updates.} Represent a user's standing information requirements as questions. The agent will identify which new evidence changes the answer, request clarification of consequential ambiguities, and prioritize updates under time or bandwidth limits. Distributed nodes provide a demanding test case, but are not the only setting.
\end{itemize}
\head{Evaluation and expected outcomes.} MEQA and MAVEN-ERE will support event-based reasoning; synthetic correction streams built from multi-document scenarios will test knowledge updates. We will measure answer accuracy, update robustness, unsupported conclusions, and cost of finding consequential changes. Baselines include standard retrieval-augmented generation and event retrieval without update tracking.
\par\head{Naval relevance.} During maritime operations, an afloat staff may need to decide whether a port is usable while delayed reports and corrections arrive from ships and shore. A decision-support agent could identify which new report changes that assessment, cite its evidence, and flag unresolved conflicts. This supports command decision making and command-and-control interests in knowledge synthesis, a common picture, and reasoning from uncertain information.
\par\head{Prior work and PI qualifications.} Our QA-based event extraction (\href{https://doi.org/10.18653/v1/2020.emnlp-main.49}{\textcolor{blue}{Du and Cardie, 2020}}) and MEQA benchmark (\href{https://openreview.net/forum?id=snNuvAOQxB}{\textcolor{blue}{Li et al., 2024}}) support the text-to-event foundation. The PI holds a 2024 NSF CAREER award on consistent event graphs and led the RESIN team's DARPA KAIROS delivery-system integration as a postdoc; the team achieved top performance among participants.

\end{document}
